\documentclass[11pt]{article}

\usepackage{amsmath,amssymb}
\usepackage{xurl}
\usepackage[hidelinks]{hyperref}
\setlength{\emergencystretch}{2em}

% Shared commands for notes in docs/paper-gaps/.

\DeclareUnicodeCharacter{2080}{\ifmmode{}_{0}\else\textsubscript{0}\fi}
\DeclareUnicodeCharacter{2081}{\ifmmode{}_{1}\else\textsubscript{1}\fi}
\DeclareUnicodeCharacter{2082}{\ifmmode{}_{2}\else\textsubscript{2}\fi}
\DeclareUnicodeCharacter{2099}{\ifmmode{}_{n}\else\textsubscript{n}\fi}

\newcommand{\Lean}{\textsc{Lean}}
\ExplSyntaxOn
\NewDocumentCommand{\leanid}{m}{
  \ifmmode
    \textsf{\detokenize{#1}}
  \else
    \group_begin:
    \urlstyle{sf}
    \tl_set:Nn \l_tmpa_tl {#1}
    \tl_replace_all:Nnn \l_tmpa_tl {\_} {_}
    \exp_args:NV \path \l_tmpa_tl
    \group_end:
  \fi
}
\ExplSyntaxOff
\providecommand{\tr}{\operatorname{tr}}
\newcommand{\tnleanIssueBase}{https://github.com/LionSR/TNLean/issues/}
\newcommand{\tnleanPrBase}{https://github.com/LionSR/TNLean/pull/}
\newcommand{\tnleanDocBase}{https://sirui-lu.com/TNLean/docs/}
\newcommand{\ghissue}[1]{\href{\tnleanIssueBase#1}{GitHub issue \##1}}
\newcommand{\ghpr}[1]{\href{\tnleanPrBase#1}{PR \##1}}
\newcommand{\doclink}[1]{\href{\tnleanDocBase#1.html}{\path{#1}}}

% Dirac notation and the identity operator, as in blueprint/src/macros/common.tex.
% \providecommand keeps note-local definitions authoritative.
\usepackage{dsfont}
\providecommand{\ket}[1]{|#1\rangle}
\providecommand{\bra}[1]{\langle #1|}
\providecommand{\braket}[2]{\langle #1 | #2 \rangle}
\providecommand{\ketbra}[2]{|#1\rangle\!\langle #2|}
\providecommand{\Id}{\mathds{1}}
\providecommand{\one}{\mathds{1}}
\providecommand{\C}{\mathbb{C}}
\providecommand{\MN}[1]{M_{#1}(\C)}
\providecommand{\MD}{M_D(\C)}

% External referencing for paper sources.  The build helper writes sanitized
% auxiliary files under build/paper-aux/ and excludes duplicated source labels.
\usepackage{xr}

% Import a generated paper auxiliary file with a caller-supplied label prefix.
% The candidate paths support builds from docs/paper-gaps/, the repository root,
% and build/paper-aux/ itself.
\newcommand{\importpaperaux}[2]{%
  \IfFileExists{../../build/paper-aux/#2.aux}{%
    \externaldocument[#1]{../../build/paper-aux/#2}%
  }{%
    \IfFileExists{build/paper-aux/#2.aux}{%
      \externaldocument[#1]{build/paper-aux/#2}%
    }{%
      \IfFileExists{#2.aux}{%
        \externaldocument[#1]{#2}%
      }{}%
    }%
  }%
}

% General external reference macros
\newcommand{\paperref}[2]{\ref{#1-#2}}
\newcommand{\papereqref}[2]{(\ref{#1-#2})}

% CPSV16 source (1606.00608 / MPDO-22-12-17-2.tex) external document and shortcuts
\importpaperaux{cpsv16-}{1606.00608/MPDO-22-12-17-2-tnlean}
\newcommand{\cpsvref}[1]{\paperref{cpsv16}{#1}}
\newcommand{\cpsveqref}[1]{\papereqref{cpsv16}{#1}}

% Verdict marker: \gapnote{<kind>}{<status>}. Kind is the policy
% classification (clarification, local-correction, scope-restriction,
% unfaithful, false-source, open-gap); status is open, wip (elimination
% underway), resolved, or historical. Machine-read by the site index;
% prints a verdict line.
\newcommand{\gapnote}[2]{\par\noindent\textbf{Verdict:} #1 (#2).\par}


\title{Dependency Audit for the Non-periodic MPS Fundamental Theorem Blueprint}
\author{}
\date{2026-05-08}

\begin{document}
\maketitle

\section{Scope}

This note records the source-faithful answer to the organizational questions in
\href{https://github.com/LionSR/TNLean/issues/1530}{issue~1530}.  It concerns
the non-periodic Fundamental Theorem route based on
\cite{Cirac2016MPDO_arXiv,Cirac2021Matrix}, not the direct periodic theorem of
\cite{DeLasCuevas2017Irreducible} and not the fixed-length algebraic theorem of
\cite{Molnar2018NormalPEPS}.  It also records which assumptions in the current
conditional theorem surface are source inputs that still have to be produced,
rather than conclusions already proved in the blueprint.

The guiding source passages are the BNT definition and characterization, the
block-injective canonical-form input, the proportional Fundamental Theorem, and
the equal-MPV corollary in \cite[lines~271--352]{Cirac2016MPDO_arXiv}, together
with the proof appendix \cite[lines~1155--1192]{Cirac2016MPDO_arXiv}.  The
review version is \cite[Section~IV]{Cirac2021Matrix}.  This note follows the
same convention as the previous paper-gap notes: the source theorem is the
standard of comparison, while formal helper structures are acceptable only when
they are identified as helper structures.

\section{Block permutation and separation}

The blueprint chapter titled ``Block Permutation and Separation'' contains more
than one mathematical ingredient.  The product-algebra part proves that an
automorphism of
\[
  \prod_{j=1}^g M_{D_j}(\mathbb C)
\]
permutes the simple factors and is inner on each factor.  This is the algebraic
mechanism used by same-structure and block-injective comparison statements.
The invariant-projection material belongs to canonical-form reduction.  The
block-separation lemmas prove that, under strict canonical-form hypotheses, a
same-MPV identity for two block-diagonal tensors separates into same-MPV
identities for the individual blocks.

For the current non-periodic Fundamental Theorem surface, this chapter is not
used wholesale.  Its visible downstream uses are the canonical-form condition
for BNT data and the same-structure comparison lemma.  The latter is the special
case in which the two tensors already have the same number of blocks, the same
bond dimensions, and the same weights; it is not the CPSV BNT matching step.  In
the source proof, the general non-periodic theorem instead proceeds through BNT
representatives, eventual linear independence, block-injective fixed-length
span, and scalar power sums.

Thus the block-permutation chapter should stay before the non-periodic theorem,
but the blueprint should not suggest that all of its material is a required
input for the remaining CPSV comparison.  The required source input still
missing from the general route is the homogeneous block-injective direct-sum
span after blocking,
\[
  \operatorname{span}\{\tilde A^w: |w|=L\}
  = \bigoplus_{j=1}^g M_{D_j}(\mathbb C),
\]
or its trace-dual separation form.  This is the input recorded separately in
\path{docs/paper-gaps/nonperiodic_mps_bnt_comparison_inputs.tex} and
\path{docs/paper-gaps/david2006_direct_sum_input.tex}.

\section{Multiplicative domain versus operator convexity}

The Kadison--Schwarz and multiplicative-domain part of the Schwarz chapter is a
real dependency of the aperiodic route.  It is used to turn peripheral
eigenvectors of normalized completely positive maps into multiplicative-domain
relations, and then into the commutation and phase-gauge rigidity needed in the
spectral and canonical-form chapters.  Consequently this material should remain
before Perron--Frobenius theory and the non-periodic Fundamental Theorem.

The operator convexity, Jensen, and Lieb-concavity material is different.  A
label audit finds no current use of that section in the non-periodic Fundamental
Theorem chapter or its BNT comparison chain.  Its natural downstream uses are
entropy, MPDO, and parent-Hamiltonian arguments, especially the open
operator-convexity and POVM-resolution work around the parent-Hamiltonian gap.
Therefore the safe first step is not to move the whole Schwarz chapter.  The
safe split, if desired in a later PR, is to keep Kadison--Schwarz and the
multiplicative domain upstream, and to move only the operator-convexity/Jensen
section to a downstream analytic chapter after the Fundamental Theorem.

\section{Which assumptions are extra}

In the conditional after-blocking theorem surface, the following data are not
extra mathematical assumptions relative to the source theorem once the upstream
canonical-form reduction has been proved: a common positive blocking length,
zero-tail equality
\[
  A^{[p]}\sim \mathbf{0}_{d^p,z_A}\oplus A_{\mathrm{nz}},
  \qquad
  B^{[p]}\sim \mathbf{0}_{d^p,z_B}\oplus B_{\mathrm{nz}},
  \qquad z_A=z_B,
\]
and trace-preserving, primitive, irreducible nonzero sectors.  These are the
formal names for the canonical-form and period-removal output in the source
route.

The remaining hypotheses are genuine proof obligations at the current theorem
surface.  They are one-site injectivity at the chosen blocking length, the BNT
representative comparison, and the proportional-decomposition data which match
representatives by a permutation, gauges, and nonzero phases,
\[
  B_{\pi(j)}^i=\zeta_j X_j A_j^i X_j^{-1}.
\]
The formal statements which assume these data do not prove them from equality
of MPV families.  They only transport them through the common blocking and then
recover the sector weights.  This is why the blueprint and the paper-gap notes
must continue to mark the passage as conditional.  In particular, this audit
does not claim that the proof gaps tracked by issues~1499, 1500, or 1501 have
been solved.

\section{Staged blueprint organization}

The chapter order can be improved without changing theorem statements.  A safe
organization is to keep the aperiodic proof dependencies first, ending with the
proof of the Fundamental Theorem.  The direct periodic theorem, symmetries and
string order, and the algebraic fixed-length theorem then form the first
post-FT block.  Entropy and MPDO material should be placed before parent
Hamiltonians, since the parent-Hamiltonian chapter depends conceptually on
state and density-operator technology rather than conversely.

The larger split of the Schwarz chapter should be staged separately.  It is a
pure organization change once the operator-convexity section is isolated, but it
touches many labels and should not be combined with proof-gap claims.

\bibliographystyle{alpha}
\bibliography{references}

\end{document}
